\documentclass[aspectratio=169]{beamer}
\usetheme{metropolis}
\usepackage{minted}
\usepackage{etoolbox}
\usepackage{booktabs}
\usepackage{tikz}
\usetikzlibrary{arrows.meta, positioning, shapes.geometric, fit}

\setminted{fontsize=\footnotesize, breaklines=true}

% Add spacing between block title and body globally (patch metropolis \nointerlineskip)
\makeatletter
\patchcmd{\metropolis@block}{\nointerlineskip}{\vskip0.4em}{}{}
\makeatother

\title{APPy: Annotated Parallelism for Python on GPUs}
\subtitle{Work in Progress}
\author{}
\date{}

% ---------------------------------------------------------------
\begin{document}
\maketitle

% ---------------------------------------------------------------
\section{Motivation}

\begin{frame}{Why GPU Acceleration Matters}
  Scientific Python workloads keep growing in scale and complexity:
  \begin{itemize}
    \item Simulation \& numerical computing
    \item Data analysis \& processing pipelines
    \item ML and non-ML scientific workloads
  \end{itemize}
  \vspace{1em}
  GPUs offer massive parallelism — but Python's default execution model is \textbf{sequential}.

  \vspace{1em}
  \begin{alertblock}{The core problem}
    Most Python programs \emph{could} benefit from GPU acceleration, but the path there is not easy.
  \end{alertblock}
\end{frame}

\begin{frame}{Existing Path 1: Library-Based Acceleration}
  \textbf{Examples:} CuPy, PyTorch ops, JAX ops, cuBLAS
  \vspace{0.5em}

  \begin{columns}
    \column{0.48\textwidth}
    \textbf{Pros}
    \begin{itemize}
      \item Easy to use
      \item Pre-optimized kernels
      \item Great for standard linear algebra \& tensor ops
    \end{itemize}

    \column{0.48\textwidth}
    \textbf{Cons}
    \begin{itemize}
      \item Many workloads cannot be expressed using predefined operators
      \item Irregular loops, custom control flow, stencil codes $\to$ not supported
      \item Must rewrite algorithms to fit the library's operator set
    \end{itemize}
  \end{columns}
\end{frame}

\begin{frame}{Existing Path 2: SIMT — Custom GPU Kernels}
  \textbf{Examples:} CUDA, Metal, OpenCL, Numba CUDA\\
  \textbf{Model:} SIMT — each thread executes \emph{independently} with its own control flow.
  \vspace{0.5em}

  \begin{columns}
    \column{0.48\textwidth}
    \textbf{Pros}
    \begin{itemize}
      \item Full control over execution
      \item Natural for irregular, control-heavy workloads (ray tracing, Mandelbrot)
      \item Potential for highest performance
    \end{itemize}

    \column{0.48\textwidth}
    \textbf{Cons}
    \begin{itemize}
      \item Must rewrite Python logic in a GPU kernel language
      \item Manual kernel launches \& memory management
      \item Collective ops (reductions, dot products) require manual thread synchronization, barrier insertion, and shared memory management — error-prone and hard to maintain
      \item Debugging: divergence, race conditions, shared-memory issues
    \end{itemize}
  \end{columns}
\end{frame}

\begin{frame}{Existing Path 3: SIMD-like — Triton}
  Triton's idea: write GPU kernels in a Python-like DSL that compiles to optimized GPU code.\\
  \textbf{Model:} SIMD-like — a thread \emph{block} collectively operates on a \emph{tile} of data.
  \vspace{0.5em}

  \begin{columns}
    \column{0.48\textwidth}
    \textbf{Strengths}
    \begin{itemize}
      \item High performance for dense, matrix/tensor-style workloads
      \item Minimal boilerplate vs.\ CUDA
      \item Strong auto-tuning infrastructure
    \end{itemize}

    \column{0.48\textwidth}
    \textbf{Limitations}
    \begin{itemize}
      \item Requires learning Triton's DSL — not pure Python
      \item Still kernel-centric: users manage blocks, program IDs, tiled memory loads
      \item \textbf{Awkward for SIMT-style code}: irregular control flow, per-thread divergence
            (e.g.\ Mandelbrot) does not fit the tile/block model
    \end{itemize}
  \end{columns}
\end{frame}

\begin{frame}{Two Paradigms, One Common Pain}
  \begin{columns}
    \column{0.48\textwidth}
    \textbf{SIMT} (CUDA, Metal)
    \begin{itemize}
      \item Each thread works independently
      \item Great for irregular, control-heavy code
      \item Collective ops require manual synchronization \& barriers — \emph{error-prone}
    \end{itemize}

    \column{0.48\textwidth}
    \textbf{SIMD-like} (Triton)
    \begin{itemize}
      \item Thread block operates on data tiles collectively
      \item Great for dense tensor/matrix ops
      \item \emph{Awkward} for irregular control flow
    \end{itemize}
  \end{columns}

  \vspace{1.2em}
  \begin{alertblock}{Shared pain points}
    Both require: manual kernel launches, explicit memory management,
    low-level boilerplate. Neither covers the full spectrum naturally.
  \end{alertblock}
\end{frame}

\begin{frame}{The Gap in the Landscape}
  No simple way to express GPU parallelism using \textbf{regular Python loops} with \textbf{minimal annotation}
  that covers \textbf{both SIMT-style} (independent threads) \textbf{and SIMD-style} (collective thread-group ops)
  within the same Python front end.

  \vspace{1em}
  \begin{center}
  \begin{tabular}{lccc}
    \toprule
    & \textbf{Triton} & \textbf{CUDA/Metal} & \textbf{APPy} \\
    \midrule
    SIMT-style code           & $\times$   & \checkmark      & \checkmark \\
    SIMD/collective ops       & \checkmark & error-prone$^*$ & \checkmark \\
    Auto host-device transfer & $\times$   & $\times$        & \checkmark \\
    No manual kernel launches & $\times$   & $\times$        & \checkmark \\
    \bottomrule
  \end{tabular}\\[0.5em]
  {\footnotesize $^*$Requires manual thread synchronization, barrier insertion, and shared memory management — a common source of subtle bugs.}
  \end{center}
\end{frame}

% ---------------------------------------------------------------
\section{APPy: Prior Work (CC 2024)}

\begin{frame}[fragile]{APPy v1 (CC 2024): Two-Layer Parallelism Model}
  \textbf{Key idea:} Annotate Python loops with pragmas; compiler generates GPU kernels.

  \begin{minted}{python}
  #pragma parallel for        # outer loop -> grid (1 thread block per iteration)
  for i in range(M):
      #pragma simd             # inner loop -> threads within that block
      for j in range(N):
          C[i, j] = A[i, j] + B[i, j]
  \end{minted}

  \vspace{0.5em}
  \textbf{Execution model:}
  \begin{itemize}
    \item Outer \texttt{\#pragma parallel for} $\to$ each iteration dispatches \textbf{one thread block}
    \item Inner \texttt{\#pragma simd} $\to$ threads within the block work collectively
    \item Essentially a \textbf{group/SIMD-like model} — each outer iteration = a cooperative group
  \end{itemize}

  \vspace{0.5em}
  No CUDA, no Triton DSL required — works well for collective inner-loop workloads.
\end{frame}

\begin{frame}[fragile]{Where APPy v1 Hit a Limit}
  CC 2024's two-layer model (outer $\to$ thread block, inner $\to$ SIMD) is essentially
  a \textbf{group model only}. It cannot naturally express \textbf{flat SIMT} code.

  \vspace{0.8em}
  \textbf{Example: Mandelbrot} — each pixel is fully independent, no inner loop to parallelize.

  \begin{minted}{python}
  #pragma parallel for        # each iteration needs a WHOLE thread block?
  for idx in range(width * height):
      # per-pixel work — no inner loop, no collective ops
      while abs(z) <= 2.0 and i < max_iter: ...
  \end{minted}

  \vspace{0.5em}
  \begin{alertblock}{The bottleneck}
    Flat, shader-style SIMT code does not fit the two-layer model.
    Launching a full thread block per pixel is wasteful and unnatural.
    CC 2024 APPy lacks a first-class \emph{individual} execution mode.
  \end{alertblock}
\end{frame}

% ---------------------------------------------------------------
\section{The Dual Programming Model}

\begin{frame}[fragile]{APPy v2: A Unified Dual Parallelism Model}
  \textbf{One annotation. Two execution modes. Compiler decides.}

  \vspace{1em}
  \begin{columns}
    \column{0.48\textwidth}
    \textbf{Individual Mode} \emph{(SIMT-style)}\\
    Each parallel iteration $\to$ one independent GPU thread.\\
    Natural for irregular, control-heavy workloads.\\
    \vspace{0.5em}
    \emph{``Assign each student a task. They work alone.''}

    \column{0.48\textwidth}
    \textbf{Group Mode} \emph{(SIMD-like)}\\
    Each parallel iteration $\to$ one GPU threadgroup.\\
    Threads inside the group coordinate on collective ops.\\
    \vspace{0.5em}
    \emph{``Divide students into groups. Students in a group can collaborate.''}
  \end{columns}

  \vspace{1.5em}
  \begin{center}
  Both modes use the \emph{same} user-facing annotation:\\[0.5em]
  \mintinline{python}{#pragma parallel for}
  \end{center}
\end{frame}

\begin{frame}{Individual vs.\ Group Mode: Visualized}
  \begin{columns}[T]
    \column{0.48\textwidth}
    \centering
    \textbf{Individual Mode}\\[0.5em]
    \begin{tikzpicture}[
      thread/.style={circle, draw, fill=mLightBrown!50, minimum size=0.65cm, font=\scriptsize},
      output/.style={draw, rounded corners, fill=mLightGreen!50, minimum width=0.7cm, minimum height=0.5cm, font=\scriptsize},
    ]
      \foreach \i in {1,2,3,4} {
        \node[thread] (t\i) at (0, {-(\i-1)*1.05}) {T\i};
        \node[output] (o\i) at (2.2, {-(\i-1)*1.05}) {out\i};
        \draw[->, thick] (t\i) -- (o\i);
      }
      \node[above=0.15cm of t1, font=\tiny, text=gray] {threads};
      \node[above=0.15cm of o1, font=\tiny, text=gray] {outputs};
    \end{tikzpicture}\\[0.5em]
    {\small Each thread: independent task}

    \column{0.48\textwidth}
    \centering
    \textbf{Group Mode}\\[0.5em]
    \begin{tikzpicture}[
      thread/.style={circle, draw, fill=blue!20, minimum size=0.6cm, font=\scriptsize},
      output/.style={draw, rounded corners, fill=mLightGreen!50, minimum width=0.7cm, minimum height=0.9cm, font=\scriptsize},
      grp/.style={draw, rounded corners=5pt, thick, dashed, inner sep=5pt},
    ]
      % Group 1
      \node[thread] (t1) at (0.0,  0.0)  {T1};
      \node[thread] (t2) at (0.85, 0.0)  {T2};
      \node[thread] (t3) at (0.0, -0.85) {T3};
      \node[thread] (t4) at (0.85,-0.85) {T4};
      \node[grp, fit=(t1)(t2)(t3)(t4), label=above:{\tiny group 1}] (g1) {};
      \node[output] (o1) at (2.6, -0.4) {out1};
      \draw[->, thick] (g1.east) -- (o1.west);
      \draw[<->, dotted, gray] (t1)--(t2); \draw[<->, dotted, gray] (t3)--(t4);
      \draw[<->, dotted, gray] (t1)--(t3); \draw[<->, dotted, gray] (t2)--(t4);

      % Group 2
      \node[thread] (t5) at (0.0, -1.9)  {T5};
      \node[thread] (t6) at (0.85,-1.9)  {T6};
      \node[thread] (t7) at (0.0, -2.75) {T7};
      \node[thread] (t8) at (0.85,-2.75) {T8};
      \node[grp, fit=(t5)(t6)(t7)(t8), label=above:{\tiny group 2}] (g2) {};
      \node[output] (o2) at (2.6, -2.3) {out2};
      \draw[->, thick] (g2.east) -- (o2.west);
      \draw[<->, dotted, gray] (t5)--(t6); \draw[<->, dotted, gray] (t7)--(t8);
      \draw[<->, dotted, gray] (t5)--(t7); \draw[<->, dotted, gray] (t6)--(t8);
    \end{tikzpicture}\\[0.5em]
    {\small Threads cooperate within a group}
  \end{columns}
\end{frame}

\begin{frame}{Individual Mode: Concept}
  \begin{block}{Analogy}
    Assign each student an independent task. They work completely independently —
    no communication, no coordination. Each one just does their own thing.
  \end{block}

  \vspace{1em}
  On the GPU: each iteration of the \texttt{\#pragma parallel for} loop maps to
  \textbf{exactly one GPU thread}.

  \vspace{1em}
  Suitable for:
  \begin{itemize}
    \item Embarrassingly parallel workloads: Mandelbrot, ray marching, prime counting
    \item Workloads where iterations are fully independent
    \item At most a \emph{final} global reduction across all iterations
  \end{itemize}
\end{frame}

\begin{frame}[fragile]{Individual Mode: Mandelbrot Example}
  \begin{minted}{python}
  @appy.jit
  def mandelbrot(out, width, height, max_iter):
      #pragma parallel for
      for idx in range(width * height):
          x = idx % width
          y = idx // width
          c = complex(x / width * 3.5 - 2.5,
                      y / height * 2.0 - 1.0)
          z, i = 0+0j, 0
          while abs(z) <= 2.0 and i < max_iter:
              z = z * z + c
              i += 1
          out[y, x] = i
  \end{minted}
  Each pixel is computed by one independent thread.
  No communication between iterations.
\end{frame}

\begin{frame}[fragile]{Individual Mode: Prime Counting Example}
  \begin{minted}{python}
  @appy.jit
  def count_primes(count, N):
      #pragma parallel for
      for n in range(2, N):
          is_prime = True
          for d in range(2, n):
              if n % d == 0:
                  is_prime = False
                  break
          if is_prime:
              count[0] += 1   # final reduction: one atomic add per thread
  \end{minted}
  Each thread independently checks one number — irregular control flow handled naturally.
  \texttt{count[0] += 1} is an atomic global reduction, efficient here because
  \textbf{primes are rare} so few threads contend on the atomic add.
\end{frame}

\begin{frame}{Individual Mode: When to Upgrade to Group Mode}
  \begin{alertblock}{Limitation of per-thread atomic reductions}
    If the reduction is \emph{unconditional} (every thread contributes),
    $N$ concurrent atomic adds cause high memory contention — inefficient.
  \end{alertblock}

  \vspace{0.8em}
  \textbf{Solution: upgrade to group mode.}
  \begin{itemize}
    \item Each threadgroup first reduces locally (fast, no global traffic)
    \item Then one atomic add per group to the global counter
    \item $\lceil N / \texttt{GROUP\_SIZE} \rceil$ atomic ops instead of $N$
  \end{itemize}

  \vspace{0.5em}
  Same \texttt{\#pragma parallel for} annotation — just rewrite the body
  to use an array expression, and the compiler handles the rest.
\end{frame}

\begin{frame}{Group Mode: Concept}
  Each \texttt{\#pragma parallel for} iteration maps to \textbf{one GPU threadgroup}.
  Threads in a group share a whiteboard (shared memory) and coordinate via barriers.

  \vspace{0.6em}
  The underlying mechanism is a \textbf{worksharing loop} (\texttt{\#pragma workshare}):
  inner loop iterations are divided among threads; implicit barrier + reduction at the end.

  \vspace{0.5em}
  Two user-facing styles — both desugar to the same worksharing loop:
  \begin{itemize}
    \item \textbf{Array expressions} — compiler rewrites to workshare loops automatically;
          concise for uniform, branchless ops
    \item \textbf{Explicit \texttt{\#pragma workshare}} — write the loop directly;
          natural for irregular control flow
  \end{itemize}
\end{frame}

\begin{frame}{Group Mode: Memory Model}
  Variables \textbf{defined inside} a group-mode body are placed in registers or shared memory
  based on one rule:

  \vspace{0.8em}
  \begin{columns}[T]
    \column{0.48\textwidth}
    \textbf{Registers} — if \emph{never} subscripted
    \begin{itemize}
      \item Private to each thread
      \item Scalars, accumulators, loop indices
    \end{itemize}

    \column{0.48\textwidth}
    \textbf{Shared memory} — if \emph{ever} subscripted
    \begin{itemize}
      \item Visible to all threads in the group
      \item Tiles, halos, local buffers
    \end{itemize}
  \end{columns}

  \vspace{0.6em}
  \textbf{Rationale:} threads cannot access each other's registers —
  shared memory is the \emph{only} inter-thread channel within a group.
  Loading into a shared tile first means global-memory latency is paid
  once per tile, not once per thread.

  \vspace{0.4em}
  \small\textit{Register overflow: Metal compiler handles spilling — APPy never force-promotes.}
\end{frame}

\begin{frame}[fragile]{Group Mode: Memory Model — 1D Jacobi}
  \textbf{1D Jacobi stencil} — each output point averages its two neighbours.

  \vspace{0.5em}
  \begin{minted}{python}
  #pragma parallel for
  for i in range(0, n, K):          # one threadgroup per tile
      buf = np.zeros(K + 2)         # subscripted  →  shared memory
      acc = 0.0                     # never subscripted  →  register

      buf[1:K+1] = a[i:i+K]        # cooperative tile load
      buf[0]     = a[i - 1]        # left halo
      buf[K+1]   = a[i + K]        # right halo
      # implicit barrier — buf fully visible to all threads
      b[i:i+K] = (buf[:-2] + buf[1:-1] + buf[2:]) / 3.0
  \end{minted}

  \vspace{0.5em}
  \texttt{buf} is in shared memory because thread \texttt{k} reads
  \texttt{buf[k-1]} and \texttt{buf[k+1]} written by neighbouring threads.
  \texttt{acc} is never shared — each thread keeps its own register copy.
\end{frame}

\begin{frame}[fragile]{Group Mode: Worksharing}
  The underlying mechanism in group mode is a \textbf{worksharing loop}:
  the inner loop is divided among threads, with an implicit barrier and reduction at the end.

  \vspace{0.5em}
  Array expressions are \textbf{syntactic sugar} — the compiler desugars them into
  explicit \texttt{\#pragma workshare} loops before codegen.

  \vspace{0.5em}
  Users can also write \texttt{\#pragma workshare} directly when the operation
  involves \textbf{irregular control flow} that is unnatural to express as array ops:

  \begin{itemize}
    \item Conditional accumulation with complex branching
    \item Multi-pass algorithms where later passes depend on earlier results
    \item Variable-length inner ranges (e.g.\ sparse row iteration)
  \end{itemize}
\end{frame}

\begin{frame}[fragile]{Group Mode: Worksharing — Radial Binning}
  \textbf{Radial binning} — average data points falling in each radial bin.
  \begin{minted}{python}
  #pragma parallel for          # one threadgroup per bin i
  for i in range(npt):
      r1, r2 = rmax*i/npt, rmax*(i+1)/npt
      s, count = 0.0, 0
      #pragma workshare         # j-loop divided among threads
      for j in range(radius.shape[0]):
          if r1 <= radius[j] < r2:
              s += data[j]; count += 1
      # implicit barrier + reduction of s, count
      res[i] = s / count
  \end{minted}

  \begin{itemize}
    \item Each thread handles a stride of the \texttt{j}-loop
    \item The \texttt{if} branch is natural — no masking gymnastics needed
    \item \texttt{s}, \texttt{count} stay in registers; reduced at the implicit barrier
  \end{itemize}
\end{frame}

\begin{frame}[fragile]{Group Mode: Symmetric Matrix Multiply (SYMM)}
  \begin{minted}{python}
  @appy.jit
  def kernel_appy(alpha, beta, C, A, B):
      C *= beta
      for i in range(C.shape[0]):
          #pragma parallel for          <- launch one group per i
          for j in range(C.shape[1]):
              C[:i, j] += alpha * B[i, j] * A[i, :i]  # collective write
              temp2[j]  = B[:i, j] @ A[i, :i]          # dot product
          C[i, :] += alpha * B[i, :] * A[i, i] + alpha * temp2
      return C
  \end{minted}

  \vspace{0.5em}
  Thread $j$ handles column $j$ — but within the group all threads cooperate on
  the slice operations and dot product via shared memory and threadgroup reductions.
\end{frame}

\begin{frame}[fragile]{Group Mode: TRMM (Triangular Matrix Multiply)}
  \begin{minted}{python}
  @appy.jit
  def kernel_appy(alpha, A, B):
      M, N = B.shape
      for i in range(M):
          #pragma parallel for
          for j in range(N):
              B[i, j] += A[i+1:, i] @ B[i+1:, j]  # dot product per (i,j)
      B *= alpha
      return B
  \end{minted}

  \vspace{0.5em}
  The \texttt{@} (dot product) inside the pfor body is a collective operation —
  the group of threads cooperates to reduce over the slice dimension.
\end{frame}

\begin{frame}[fragile]{Individual vs.\ Group Mode: Side by Side}
  \begin{center}
  \small
  \begin{tabular}{lll}
    \toprule
    & \textbf{Individual} & \textbf{Group} \\
    \midrule
    pfor iter maps to   & 1 thread          & 1 threadgroup \\
    Thread coordination & None              & Shared memory + barriers \\
    Inner mechanism     & ---               & \texttt{\#pragma workshare} loop \\
    Inner notation      & ---               & array exprs (sugar) or explicit workshare \\
    Control flow        & Per-thread        & Arbitrary \\
    Barrier at end      & No                & Implicit \\
    Suitable for        & Independent tasks & Collective ops per iteration \\
    Examples            & Mandelbrot        & TRMM, SYMM, radial binning \\
    \midrule
    Annotation          & \multicolumn{2}{c}{\texttt{\#pragma parallel for}} \\
    \bottomrule
  \end{tabular}
  \end{center}
\end{frame}

% ---------------------------------------------------------------
\section{Compiler Design}

\begin{frame}{Compiler Pipeline}
  \begin{center}
  \scalebox{0.56}{
  \begin{tikzpicture}[
    stage/.style={draw, rounded corners, minimum width=4.2cm, minimum height=0.75cm,
                  font=\small, fill=mLightBrown!25, align=center},
    decision/.style={draw, diamond, aspect=2.5, minimum width=4.2cm, minimum height=0.75cm,
                     font=\small, fill=yellow!25, align=center},
    codegen/.style={draw, rounded corners, minimum width=3.5cm, minimum height=0.75cm,
                    font=\small, align=center},
    arr/.style={->, thick},
  ]
    \node[stage] (src) {Python source + \texttt{\#pragma parallel for}};
    \node[stage, below=0.5cm of src] (fe) {Frontend: parse \& normalize AST};
    \node[decision, below=0.5cm of fe] (detect) {Array exprs in pfor body?};

    \node[codegen, below left=0.8cm and 1.2cm of detect, fill=mLightBrown!40]
      (indiv) {\textbf{Individual mode}\\\small 1 thread / iteration};
    \node[codegen, below right=0.8cm and 1.2cm of detect, fill=blue!15]
      (group) {\textbf{Group mode}\\\small 1 threadgroup / iteration};

    \node[stage, below=0.7cm of group, fill=yellow!20, font=\small]
      (lower) {Lower array exprs $\to$ \texttt{\#pragma workshare} loops};
    \node[codegen, below=0.7cm of lower, fill=blue!25, font=\small]
      (wsmode) {Worksharing codegen\\\small (divide loop, barrier, reduce)};

    \node[stage, below=4.5cm of detect] (backend) {Backend codegen (Metal / Triton)};
    \node[stage, below=0.5cm of backend, fill=mLightGreen!40] (out) {GPU kernel};

    \draw[arr] (src)    -- (fe);
    \draw[arr] (fe)     -- (detect);
    \draw[arr] (detect) -- node[left, font=\scriptsize]  {no}  (indiv);
    \draw[arr] (detect) -- node[right, font=\scriptsize] {yes} (group);
    \draw[arr] (group)  -- (lower);
    \draw[arr] (lower)  -- (wsmode);
    \draw[arr] (indiv)  |- (backend);
    \draw[arr] (wsmode) |- (backend);
    \draw[arr] (backend) -- (out);
  \end{tikzpicture}
  }
  \end{center}
\end{frame}

% ---------------------------------------------------------------
\section{Supported Workloads \& Examples}

\begin{frame}{NPBench Kernels Supported}
  \begin{center}
  \begin{tabular}{llll}
    \toprule
    \textbf{Kernel} & \textbf{Mode} & \textbf{Key operation} & \textbf{Annotation} \\
    \midrule
    SAXPY / vec-add & Individual & Pointwise & \texttt{\#pragma parallel for} \\
    GELU            & Individual & Pointwise & \texttt{\#pragma parallel for} \\
    Mandelbrot      & Individual & Irregular & \texttt{\#pragma parallel for} \\
    GESUMMV         & Group & \texttt{np.sum} & \texttt{\#pragma parallel for} \\
    GEMVER          & Group & \texttt{@} (dot) & \texttt{\#pragma parallel for} \\
    TRMM            & Group & \texttt{@} (dot) & \texttt{\#pragma parallel for} \\
    SYMM            & Group & \texttt{@} + scatter & \texttt{\#pragma parallel for} \\
    Gramschmidt     & Group & \texttt{@} (dot) & \texttt{\#pragma parallel for} \\
    SYRK / SYR2K    & Group & Slice accumulate & \texttt{\#pragma parallel for} \\
    Jacobi-1D/2D    & Group & Stencil & \texttt{\#pragma parallel for} \\
    \bottomrule
  \end{tabular}
  \end{center}
\end{frame}

\begin{frame}[fragile]{Writing APPy vs.\ NumPy Reference}
  \textbf{NumPy reference (TRMM):}
  \begin{minted}{python}
  for i in range(B.shape[0]):
      for j in range(B.shape[1]):
          B[i, j] += np.dot(A[i+1:, i], B[i+1:, j])
  B *= alpha
  \end{minted}

  \textbf{APPy kernel:}
  \begin{minted}{python}
  @appy.jit
  def kernel(alpha, A, B):
      for i in range(B.shape[0]):
          #pragma parallel for
          for j in range(B.shape[1]):
              B[i, j] += A[i+1:, i] @ B[i+1:, j]
      B *= alpha
  \end{minted}

  \vspace{0.5em}
  The APPy version is essentially the NumPy version with one added pragma.
\end{frame}

% ---------------------------------------------------------------
\section{Progressive Optimization Path}

\begin{frame}[fragile]{Individual $\to$ Group: A Natural Upgrade Path}
  A key benefit of APPy's unified model: users can \textbf{start simple} and
  \textbf{upgrade incrementally} — using the same annotation.

  \vspace{0.8em}
  \textbf{Example: Stencil computation}

  \textbf{Step 1 — Individual mode} (simple, correct):
  \begin{minted}{python}
  #pragma parallel for
  for idx in range(N * N):   # each thread computes one output element
      i, j = idx // N, idx % N
      out[i, j] = 0.25 * (A[i-1,j] + A[i+1,j] + A[i,j-1] + A[i,j+1])
  \end{minted}

  \textbf{Step 2 — Group mode} (tiled, better data locality):
  \begin{minted}{python}
  #pragma parallel for        # each group handles one tile row
  for i in range(0, N, TILE):
      out[i:i+TILE, :] = 0.25 * (A[i-1:i+TILE-1, :] + A[i+1:i+TILE+1, :]
                                + A[i:i+TILE, :-1]   + A[i:i+TILE, 1:])
  \end{minted}

  Same pragma, same Python front end — performance can be improved by rewriting
  to use array expressions, without leaving APPy.
\end{frame}

% ---------------------------------------------------------------
\section{Preliminary Results}

\begin{frame}{Benchmark Setup}
  \begin{itemize}
    \item \textbf{Platform:} Apple M3 (MacBook Air) — Metal GPU backend
    \item \textbf{Benchmarks:} NPBench kernels (PolyBench + custom) — all float32
    \item \textbf{Baselines:} NumPy, Numba (best mode per kernel)
    \item \textbf{Metric:} Median wall-clock time (ms), 5 runs
    \item APPy uses \textbf{a single \texttt{\#pragma parallel for}} per kernel
  \end{itemize}
\end{frame}

\begin{frame}{Results: Individual Mode — Irregular Workloads (\emph{paper} size)}
  \centering
  \footnotesize
  \begin{tabular}{lrrrl}
    \toprule
    Kernel & NumPy & Numba & APPy Metal & vs.\ NumPy \\
    \midrule
    \texttt{mandelbrot} (1000$\times$1000, iter=200) & 1407 ms & 239 ms & \textbf{3 ms} & \textbf{469$\times$} \\
    \texttt{prime\_counting} ($N{=}10^6$)            & 2274 ms & --- & \textbf{21 ms} & \textbf{108$\times$} \\
    \bottomrule
  \end{tabular}

  \vspace{0.8em}
  \begin{block}{Why such large speedups?}
    Each iteration is \emph{fully independent} with \emph{irregular control flow}
    (early exit, variable work per thread). NumPy/Numba cannot vectorise these efficiently.
    APPy launches one GPU thread per pixel/number — pure individual mode.
  \end{block}
\end{frame}

\begin{frame}{Results: Group Mode — Dense Linear Algebra (\emph{paper} size)}
  \centering
  \footnotesize
  \begin{tabular}{lrrrl}
    \toprule
    Kernel & NumPy & Numba & APPy Metal & vs.\ NumPy \\
    \midrule
    \texttt{syr2k}       & 3968 ms & 975 ms  &  64 ms & \textbf{62$\times$} \\
    \texttt{syrk}        & 1458 ms & 540 ms  &  31 ms & \textbf{47$\times$} \\
    \texttt{gemver}      &  237 ms & 141 ms  &  31 ms & \textbf{7.6$\times$} \\
    \texttt{gesummv}     &  249 ms & 117 ms  & 158 ms & 1.6$\times$ \\
    \texttt{symm}        & 4448 ms & 673 ms  & 2064 ms & 2.2$\times$ \\
    \texttt{trmm}        & 1605 ms & 292 ms  & 1309 ms & 1.2$\times$ \\
    \texttt{gramschmidt} &   75 ms &  26 ms  &  155 ms & 0.5$\times$ \\
    \bottomrule
  \end{tabular}

  \vspace{0.6em}
  \begin{block}{Takeaway}
    Memory-bandwidth-bound kernels (\texttt{syrk}, \texttt{syr2k}, \texttt{gemver}) see
    large GPU speedups. Loop-heavy kernels with sequential data dependencies are harder to parallelise.
  \end{block}
\end{frame}

\begin{frame}{Results: Stencil Kernels — Dispatch Overhead (\emph{M} size)}
  \centering
  \footnotesize
  \begin{tabular}{lrrrl}
    \toprule
    Kernel & NumPy & Numba & APPy Metal & vs.\ NumPy \\
    \midrule
    \texttt{fdtd\_2d} \emph{(M)}        &  34 ms &  61 ms &  113 ms & 0.30$\times$ \\
    \texttt{jacobi\_2d} \emph{(M)}      &  45 ms &  30 ms &  116 ms & 0.39$\times$ \\
    \texttt{floyd\_warshall} \emph{(M)} &  41 ms &  38 ms &  253 ms & 0.16$\times$ \\
    \texttt{floyd\_warshall} \emph{(paper, $N{=}2800$)} & 17527 ms & 3140 ms & 4461 ms & \textbf{3.9$\times$} \\
    \texttt{jacobi\_1d} \emph{(M)}      &  60 ms & 110 ms & 10704 ms & 0.006$\times$ \\
    \bottomrule
  \end{tabular}

  \vspace{0.6em}
  \begin{alertblock}{Bottleneck: per-kernel dispatch overhead (${\approx}240\,\mu$s/dispatch)}
    Stencil time-step loops require one GPU dispatch per step (no global sync in Metal).
    \texttt{jacobi\_1d}: 4000 steps $\times$ 2 dispatches $= 8000$ launches.
    Measured: ${\approx}165\,\mu$s Metal overhead $+$ ${\approx}85\,\mu$s actual GPU compute per step.
  \end{alertblock}
\end{frame}

% ---------------------------------------------------------------
\section{Discussion \& Future Work}

\begin{frame}{What We've Shown So Far}
  \begin{itemize}
    \item A unified programming model that supports two GPU parallelism paradigms
          with a \textbf{single annotation}
    \item A compiler that automatically infers the mode and generates correct,
          efficient GPU code
    \item Support for compound reduction expressions (\texttt{+=\ np.sum(\ldots)},
          \texttt{@} inside expressions)
    \item Backends: \textbf{Metal} (macOS) and \textbf{Triton} (NVIDIA/Linux)
    \item Validated on a range of NPBench kernels
  \end{itemize}
\end{frame}

\begin{frame}{Limitations \& Open Questions}
  \begin{itemize}
    \item \textbf{1D array expressions only} inside pfor bodies — no 2D slice ops yet
    \item \textbf{No auto-tuning} for threadgroup sizes (currently fixed)
    \item \textbf{\texttt{\#pragma simd}} is an internal annotation that could be
          renamed \texttt{\#pragma group} to better reflect the semantics
    \item How to handle \textbf{memory allocation} for intermediate arrays inside pfor bodies?
    \item Formal correctness argument for the guard pass is informal today
  \end{itemize}
\end{frame}

\begin{frame}{Future Work}
  \begin{itemize}
    \item \textbf{Rename \texttt{\#pragma simd} $\to$ \texttt{\#pragma group}}:
          SIMD is too narrow — group captures that threads \emph{cooperate},
          not just execute in lockstep
    \item \textbf{2D and higher-dimensional array expressions} inside pfor loops
    \item \textbf{Common backend passes}: consolidate shared logic (e.g.\
          \texttt{ExtractReductionSubexprs}) rather than duplicating per backend
    \item \textbf{Performance evaluation} on NPBench against NumPy, CuPy, and Triton baselines
    \item \textbf{Auto-tuning} for threadgroup size and inner loop tile size
    \item \textbf{Worksharing sub-mode} (\texttt{\#pragma workshare}):
          implement compiler support for dividing an inner loop among threads
          in a group, with implicit barrier and reduction at loop end.
          Enables algorithms with irregular per-element control flow that are
          awkward to express as array expressions — e.g.\ \textbf{bitonic sort}
          (cooperative in-place sort of a shared-memory tile across multiple
          compare-and-swap passes) and \textbf{sparse matrix-vector multiply}
          (variable-length row iteration in CSR format)
  \end{itemize}
\end{frame}

\begin{frame}[fragile]{Summary}
  \begin{block}{APPy's key contribution}
    A Python GPU programming model with a \textbf{single pragma} that supports
    two complementary parallelism paradigms — individual and group mode —
    without requiring users to think about threads, blocks, or shared memory.
  \end{block}

  \vspace{1em}
  \textbf{Individual mode:} ``Each student works alone on their task.''\\
  \textbf{Group mode:} ``Students form groups and collaborate.''\\[0.5em]
  Both expressed with \mintinline{python}{#pragma parallel for}.

  \vspace{1em}
  \begin{center}
    \large Questions \& Feedback Welcome
  \end{center}
\end{frame}

\end{document}
